\section{Hexagonal Architecture \& DIP}

The fundamental challenge in designing a distributed machine learning pipeline is preventing software entropy. If the algorithmic logic (e.g., semantic vector distance) becomes deeply coupled to the I/O mechanisms (e.g., PostgreSQL connections, FastAPI HTTP requests), the system becomes rigid, impossible to unit test deterministically, and highly susceptible to cascading failures.

To enforce mathematical purity, the Semantic Document Engine strictly adheres to \textbf{Hexagonal Architecture} (Ports and Adapters), governed by the \textbf{Dependency Inversion Principle (DIP)}.

\begin{semanticpanel}[The Dependency Inversion Principle]
High-level modules must not depend on low-level modules. Both must depend on abstractions. Furthermore, abstractions must not depend on details; details must depend on abstractions.
\end{semanticpanel}


In practice, this means the dependency arrows must unconditionally point \textit{inward} toward the Core Domain, never outward. We achieve this inversion of control through two distinct structural boundaries:

\begin{itemize}
    \item \textbf{Ports (The Interfaces):} The Core Domain dictates the contracts it requires to operate. It does not care how the data is retrieved, only that the data conforms to a specific shape. In Python, these Ports are implemented as Abstract Base Classes (\texttt{abc.ABC}) or structural \texttt{typing.Protocol} classes.
    \item \textbf{Adapters (The Implementations):} The Infrastructure layer contains the concrete code that interacts with the outside world (e.g., the network, the disk). An adapter class imports the Port from the Core Domain and physically implements its methods.
\end{itemize}

To illustrate this, consider how our system fetches a vector. Instead of the Core Domain importing a PostgreSQL driver, the Core defines a \texttt{DocumentRepository} protocol. The Infrastructure layer then provides a \texttt{PgVectorRepository} that fulfills this protocol.

\begin{lstlisting}[caption={Dependency Inversion using Python Protocols}]
from typing import Protocol

1. THE PORT (Core Domain)
class DocumentRepository(Protocol):
    def get_vector(self, doc_id: str) -> list[float]:
        ...

2. THE ADAPTER (Infra Layer)
class PgVectorRepository:
    def get_vector(self, doc_id: str) -> list[float]:
        SQL logic and network calls happen here ...
        return [0.1, 0.2, 0.3]
\end{lstlisting}

This fundamentally flips traditional application design. Instead of the Core Domain relying on the database to exist, the Core dictates the mathematical reality, and it is the database's responsibility to conform to that reality.


\vspace{1em}
\begin{figure}[ht]
    \centering
    \begin{tikzpicture}[
        node distance=1.5cm and 2cm,
        box/.style={rectangle, draw=DarkSlate, thick, rounded corners, minimum width=3.5cm, minimum height=1cm, align=center, font=\sffamily},
        core/.style={box, fill=LightPink!20, draw=LightPink, very thick},
        infra/.style={box, fill=LightBlue!20, draw=LightBlue, very thick},
        arrow/.style={-{Stealth[scale=1.2]}, thick, DarkSlate}
    ]

        % Nodes
        \node[core] (domain) {\textbf{Core Domain} \\ \footnotesize (Pure Math \& Entities)};
        \node[core, above=of domain] (interfaces) {\textbf{Core Interfaces} \\ \footnotesize (Protocols / Ports)};

        \node[infra, left=2.5cm of interfaces] (postgres) {\textbf{PostgreSQL} \\ \footnotesize (Repository Adapter)};
        \node[infra, right=of interfaces] (fastapi) {\textbf{FastAPI} \\ \footnotesize (Web Adapter)};

        \node[infra, above=2.5cm of interfaces] (celery) {\textbf{Celery} \\ \footnotesize (Worker Adapter)};

        % Background Box for Core
        \begin{pgfonlayer}{background}
            \node[draw=LightPink, fill=LightPink!5, dashed, rounded corners, fit=(domain)(interfaces), inner sep=0.5cm, label={[DarkSlate, font=\bfseries]north:Isolated System}] (corebox) {};
        \end{pgfonlayer}

        % Dependency Arrows
        \draw[arrow] (interfaces) -- (domain) node[midway, right, font=\small] {uses};
        \draw[arrow] (postgres) -- (interfaces) node[midway, above, inner xsep=0.5cm, font=\small] {implements};
        \draw[arrow] (fastapi) -- (interfaces) node[midway, above, font=\small] {uses};
        \draw[arrow] (celery) -- (interfaces) node[pos=0.2, right, font=\small] {uses};

    \end{tikzpicture}
    \caption{The Dependency Inversion Graph. All I/O adapters depend strictly on the isolated Core Interfaces. The Core Domain remains mathematically unaware of the external network topology.}
    \label{fig:hexagonal}
\end{figure}
\vspace{1em}

%By structuring the application as an inward-flowing Directed Acyclic Graph (DAG), we guarantee that a change in the SQL schema, a version bump in the FastAPI router, or a shift in the Celery configuration cannot physically alter the state-space of the Machine Learning mathematics. The core vector calculations remain entirely isolated in pure Python memory, completely unaware of whether their inputs arrived via a TCP network socket or a local unit test sandbox.

This inward-pointing structure represents a profound shift in standard \textbf{Object-Oriented Analysis and Design (OOAD)}. In a traditional \textbf{Layered Architecture}, dependencies flow top-down: the Presentation Layer depends on the Business Logic, which in turn strictly depends on the Database Layer. This approach fatally couples the application's core to its storage mechanism.

To resolve this, our system adopts the principles of \textbf{Clean Architecture}, formalized by Robert C. Martin in 2012. Clean Architecture synthesizes Hexagonal design into a topological model of concentric rings. The inner rings are reserved exclusively for pure \textit{Entities} (our vector mathematics) and \textit{Use Cases} (our semantic tracking logic). The outer rings encapsulate all interfaces, web frameworks, external APIs, and databases. The golden rule of Clean Architecture is strict dependency inversion: source code dependencies must only point inward, from the outer infrastructure rings toward the inner domain rings.

We enforce this boundary using \textbf{Inversion of Control (IoC)}. Often summarized by the Hollywood Principle---``Don't call us, we'll call you''---IoC dictates that the Core Domain never actively instantiates or calls external infrastructure. Instead, the Core Domain passively waits. In modern \textbf{Event-Driven Programming}, an external framework (such as our asynchronous \texttt{FastAPI} event loop) listens for external stimuli and dynamically triggers the Core Domain's pure functions.

To make this externalization work without breaking the system, we rely heavily on \textbf{Dependency Injection}. Rather than the Core Domain executing a command to create its own database connection, the Infrastructure layer instantiates the \texttt{PgVectorRepository} and \textit{injects} it (passes it as an argument) into the Core Domain at runtime. This physically externalizes the infrastructure, ensuring loose coupling between the mathematical logic and the disk storage.

%Ultimately, by structuring the application as an inward-flowing Directed Acyclic Graph (DAG), we guarantee that a change in the SQL schema, a version bump in the routing framework, or a shift in the \texttt{Celery} configuration cannot physically alter the state-space of the Machine Learning mathematics. The core vector calculations remain entirely isolated in pure Python memory, completely unaware of whether their inputs arrived via a TCP network socket or a local unit test sandbox.

Ultimately, the primary goal of this architecture is to structure the application's module dependencies as an inward-flowing \textbf{Directed Acyclic Graph (DAG)}. In graph theory, a DAG is a topological structure where edges point in exactly one direction (Directed) and never loop back upon themselves (Acyclic). In software engineering, this means if the Infrastructure module imports the Core Domain, the Core Domain is strictly prohibited from importing the Infrastructure module. This acyclic property is critical: it completely eliminates circular import errors---a notorious cause of runtime deadlocks in Python---and guarantees that the system can be logically reasoned about from the inside out.

By enforcing this strict, one-way topological flow, we guarantee that changes to the external environment cannot propagate inward to corrupt the mathematical logic. For instance, consider a modification to the \textbf{SQL schema}---the rigid, tabular blueprint defining the exact tables, columns, constraints, and data types stored in our PostgreSQL database. If a database administrator alters a table structure or swaps a standard array for a \texttt{pgvector} tensor, that structural mutation remains trapped entirely within the outer Repository Adapter. The Core Domain never sees the SQL query; it only sees the pure Python \texttt{list[float]} returned by the Port.

Similarly, a version bump in the \texttt{FastAPI} web router, or a shift in the \texttt{Celery} asynchronous worker configuration, cannot physically alter the state-space of the Machine Learning computations. The core vector mathematics remain completely isolated in pure Python memory. They compute topological distances and semantic drift completely unaware of whether their input tensors arrived via a live TCP network socket, a distributed message queue, or a local \texttt{pytest} sandbox.
